\section{Approach and Work Packages}

The project is organised into five interconnected work packages following a de-risked sequence. A physics-based backbone is established in Year 1, and hybrid machine learning extensions are layered on top in Years 2 and 3.

\paragraph{WP1. VIC Denmark with Groundwater Coupling.}
I will build a national VIC 5 implementation at 1 to 10\nobreakspace km, parameterised from \textbf{VICGlobal} and refined with Danish hydrostratigraphy from \textbf{GEUS} and Jupiter wells. The groundwater module will start from a lumped per-cell aquifer, following the philosophy of \textbf{WGHM} and \textbf{SIMGM} (Niu \textit{et al.}\nobreakspace 2007), with a clear escalation path to 2D lateral coupling (Zhang and Liang 2021) only if validation justifies the added complexity. This staged design directly addresses the most acute structural limitation of VIC, the leaky bucket lower boundary, without committing prematurely to a fully distributed \textbf{MODFLOW}-style solver.

\paragraph{WP2. Joint Calibration and Multi-Source Data Assimilation.}
I will implement augmented state ensemble Kalman filtering (\textbf{LETKF}) and Constrained Bayesian \textbf{MCMC} (\textbf{ConBay}) for simultaneous parameter calibration and state estimation. Observation operators will be built for \textbf{GRACE FO} terrestrial water storage (vertical column sum, sub-monthly L1B following Retegui Schiettekatte \textit{et al.}\nobreakspace 2025), \textbf{SMAP} L4 root zone and \textbf{ESA CCI} surface soil moisture (\textbf{CDF} matching, De Lannoy and Reichle 2016), routed streamflow (Lohmann routing or CaMa Flood), and groundwater heads (via effective porosity). The framework will build directly on the group's \textbf{PyGLDA} infrastructure (Yang \textit{et al.}\nobreakspace 2025, \textit{GMD} 18:6195), addressing the Denmark-to-GRACE scale mismatch through sub-monthly assimilation, ConBay merging with SMAP (Mehrnegar \textit{et al.}\nobreakspace 2023), and proper Fennoscandian \textbf{GIA} correction (\textbf{ICE-6G\_D}).

\paragraph{WP3. Hybrid Extension via Differentiable Hydrology.}
In Year 2 I will train a differentiable VIC variant, following the dPL approach of Tsai \textit{et al.}\nobreakspace (2021) and Feng \textit{et al.}\nobreakspace (2022), on \textbf{CAMELS DK} (Liu \textit{et al.}\nobreakspace 2025, \textit{ESSD} 17:1551), Jupiter piezometric records, and satellite forcing. This replaces per-basin \textbf{SCE-UA} calibration with end-to-end gradient-based parameter learning across all 3,330 basins, including ungauged ones. An \textbf{LSTM} or Transformer residual post-processor (\textbf{NeuralHydrology HybridModel}) will correct systematic biases, directly addressing the Liu \textit{et al.}\nobreakspace (2024, \textit{HESS} 28:2871) finding that pure LSTMs fail in groundwater-dependent Danish catchments.

\paragraph{WP4. Probabilistic Four-Pillar Combined Drought Index.}
For each grid cell and time step I will compute standardised sub-indices spanning the full propagation chain: atmospheric (\textbf{SPI} or \textbf{SPEI} from VIC forcing), agricultural (\textbf{SSMI} from root zone soil moisture constrained by \textbf{SMAP} and \textbf{ESA CCI}), surface hydrological (\textbf{SSI} or \textbf{SRI} from routed runoff constrained by streamflow observations), and groundwater (\textbf{SGI} from the new VIC groundwater state constrained by \textbf{GRACE FO} and Jupiter). Sub-indices will be aggregated via Gaussian or vine copulas, producing a probabilistic \textbf{MCDI} for each cell and time step with full posterior probabilities of falling within \textbf{USDM} or \textbf{CDI} severity categories. To my knowledge, a probabilistic, four-pillar, satellite data-assimilation-driven, copula-aggregated drought index with an explicit groundwater pillar has not yet been produced for Denmark. This is the unifying scientific deliverable of the thesis.

\paragraph{WP5. Short- to Seasonal-Scale Drought Forecasting.}
I will force the calibrated VIC DK with bias-corrected \textbf{SEAS5} and \textbf{AIFS ENS} or GenCast (Price \textit{et al.}\nobreakspace 2025) for 1 to 6\nobreakspace month outlooks. Evaluation will cover hindcasts from 1993 to present and the 2018, 2022, and 2023 events, using \textbf{CRPS}, \textbf{Brier score}, and \textbf{ROC}, with reverse \textbf{ESP} diagnostics to disentangle initial-condition skill from forcing skill. Benchmarks include \textbf{EFAS Seasonal} and \textbf{EDgE}.
